\documentclass[twocolumn,showpacs,preprintnumbers,amsmath,amssymb,aps,prb]{revtex4-1}

\usepackage{graphicx}
\usepackage{dcolumn}
\usepackage{bm}
\usepackage{hyperref}
\usepackage[utf8]{inputenc}
\usepackage[spanish]{babel}
\usepackage{float}

\begin{document}

\title{Simulación Monte Carlo de las Propiedades Magnéticas y Efecto Magnetocalórico en un Antiferromagneto Tipo-A: Estudio del Modelo de Ising 2D}

\author{Sofía Moscoso Ortiz, Mario César Barrera Salas}
\affiliation{Instituto de Física, Universidad de Antioquia}
\date{\today}

\begin{abstract}
En este trabajo se estudian las propiedades magnéticas y el efecto magnetocalórico (EMC) en un modelo de Ising 2D sobre una red cuadrada que representa la física esencial de un antiferromagneto tipo-A. Mediante simulaciones de Monte Carlo empleando el algoritmo de Metrópolis, se analizan el parámetro de orden magnetización de subred $M_s$, la susceptibilidad $\chi_s$, el calor específico $C_v$ y el cumulante de Binder $U$ bajo condiciones de frontera periódicas y libres para tamaños de red de $L=16$ y $L=24$. Se determina la temperatura de Néel ($T_N$) y se evalúa la variación de entropía magnética ($\Delta S_M$) y el poder de enfriamiento relativo (RCP) mediante la relación de Maxwell. Ademas, Se hace una comparacion de este modelo Ising 2D con el modelo de Heisenberg 3D. 
\end{abstract}
\maketitle


\input{01_introduction}
\input{02_resultados}
\input{03_conclusiones}
\input{04_referencias}

\clearpage
\appendix
\input{Apendice}

\end{document}
